\section{Two successful levels and upper decoding}\label{sec:upper}

We show that appreciable test acceptance produces two local bilinear-form
decoders.  Their outputs agree on the projected lower space with a probability
depending only on the upper row count.  The left output is always a normalized
product form whose restriction to its displayed lower rows has rank at most
one.  Keeping this probability independent of the lower row count will allow
us to choose that count afterward.

Fix a legal labeling of the partition vertices, extended to unused legal keys
as in Section~\ref{sec:construction}.  All responses below are evaluations
restored from this fixed labeling.  For a candidate test at level $p$, write
$\mathrm{Acc}_p$ for its acceptance event.  Sample the directions for all candidate
tests independently, using the same left query and sampled path.

\subsection{Selecting two levels and a quotient table}

Suppose that the labeling accepts with probability at least $\eps$, where
$0<\eps<1$, and choose $d\eps\ge4$.  If
$N=\sum_{p=1}^d\one_{\mathrm{Acc}_p}$, then
\[
 \E[N(N-1)]\ge (\E N)^2-\E N
       \ge d^2\eps^2-d\eps.
\]
Consequently some $i<j$ satisfies
\begin{equation}\label{eq:upper-pair}
 \Prob(\mathrm{Acc}_i\cap \mathrm{Acc}_j)
 \ge \frac{d^2\eps^2-d\eps}{d(d-1)}
 \ge \frac{3\eps^2}{4}.
\end{equation}
No independence between the acceptance events is used here.  Pass to the
modified experiment at level $i$ from Section~\ref{sec:sampling}: each child
of $D_i$ is independently projected with probability $\beta_i=n_i^{-2/3}$,
and all calls stopped at $D_i$, its own matrix, and its descendant matrices
are independent uniform draws in their new spaces.  Write $\Prob_i$ for
this law, with all functions also pulled back to the original tuple $E$.
By Lemmas~\ref{lem:cut-tv} and~\ref{lem:sparse-tv}, we can require the
comparison error in \eqref{eq:upper-pair} to be at most $\eps^2/4$.
Thus, throughout this section,
\begin{equation}\label{eq:upper-c}
 \Prob_i(\mathrm{Acc}_i\cap \mathrm{Acc}_j)\ge c,\qquad c=\eps^2/2.
\end{equation}
All eventual parameter choices will satisfy the stated bounds for every
possible pair $i<j$.

The comparison at the upper cut forgets which level-$i$ descendant lies on
the path.  We therefore build one response table using the displayed rows
at every level-$i$ descendant of $D_j$.
Put $H=H_{D_j}(E)$, $K=K_j=\F^{\ell}$, $\ell=\ell_j$, and
$M=M_{D_j}$.  The data
\[
 P=(E,\text{path through }D_j,(M_D)_{D\ne D_j})
\]
contain actual functions on the original domains.  The record $P$ does
not include an explicit tag for the selected level-$i$ descendant of $D_j$,
and does not append the projection mask or the questions in $O$.
The retained arrays may carry statistical information about these latent
choices; the law comparisons already apply to the entire retained record.
Let $\mathcal D_i(D_j)$ be the set of \emph{all} level-$i$ descendants of
$D_j$.  Write $s_{C,t}\in H_C(E)$ for row $t$ of $M_C$, where
$C\in\mathcal D_i(D_j)$ and $1\le t\le\ell_i$.  Define
\begin{equation}\label{eq:upper-B}
 B=B(P)=\Span\bigl(\{1\}\cup
       \{s_{C,t}s_{C,u}:C\in\mathcal D_i(D_j),\ 1\le t,u\le\ell_i\}\bigr)
       \subseteq H.
\end{equation}
The inclusion follows from $H_C^2\subseteq H_{D_j}$.  Fix, by deterministic
linear algebra on $P$, a linear section $s:H/B\longrightarrow H$ of the
quotient map $\pi:H\longrightarrow H/B$.  Apply both maps row by row.

For $\bar M\in K\otimes(H/B)$, query the left table with upper matrix
$s(\bar M)$ and all other matrices specified by $P$.  Retain the response
\begin{equation}\label{eq:upper-representative}
 f_P(\bar M)=\bigl(u,(y_C)_{C\in\mathcal D_i(D_j)}\bigr),
 \qquad u\in K,\quad y_C\in K_i.
\end{equation}
Thus the response contains the upper evaluation and the evaluations at every
level-$i$ descendant, without a tag marking the selected descendant.  This
finite range can be large; none of our inverse bounds depends on its size.

To recover the response at the original matrix $M$, put $\bar M=\pi(M)$ and write
$\Delta=M-s(\bar M)\in K\otimes B$.  Addition of $\Delta$ to the upper
rows is an invertible triangular transformation of the combined evaluation
tuple, since each row of $\Delta$ is a constant plus a sum of products of
unchanged lower rows.  Lemma~\ref{lem:folding} therefore shows that all
$y_C$ are unchanged and that the upper response changes from $u$ to
$u+\Delta b_y$, where $b_y$ evaluates the generators of $B$ according to
$1\mapsto1$ and $s_{C,t}s_{C,u}\mapsto (y_C)_t(y_C)_u$.
This rule is a well-defined linear functional: one legal assignment realizes
the \emph{entire} representative response, and its evaluation respects every
pointwise linear relation among the generators in \eqref{eq:upper-B}.

Define a usefulness mark by
\begin{equation}\label{eq:upper-mark}
 U(P,\bar M)=\one\{\Prob_i(\mathrm{Acc}_i\mid P,\bar M)\ge\kappa\},
 \qquad \kappa=c/2,
\end{equation}
setting it to zero off the support of the conditioning data.  The conditional
probability averages over the selected lower path, its omitted direction,
the $B$-component of $M$, and any other unexposed randomness.  Define the
partial table $\widetilde f_P$ to equal $f_P$ where $U=1$, and to be
undefined ($\bot$) elsewhere.  We fix these marks once using $\Prob_i$ and use exactly the same
partial tables under all subsequent laws.

\begin{lemma}[Defined equality from a two-response fiber]
\label{lem:useful-equality}
Under \eqref{eq:upper-c}, resampling the level-$j$ bucket indexed by its
omitted direction gives two matrices $M,M'$ such that
\[
 \Prob_i\bigl(\widetilde f_P(\bar M)
      =\widetilde f_P(\bar M')\ne\bot\bigr)\ge c^2/8.
\]
If the cut comparison for this enlarged experiment has total-variation
error at most $c^2/16$, then under the unrestricted level-$j$ cut law the
same frozen partial tables have average defined equality at least
$\theta=c^2/16$ on the ordinary uniform rank-one walk on $K\otimes(H/B)$.
\end{lemma}

\begin{proof}
The part of $\mathrm{Acc}_i\cap \mathrm{Acc}_j$ on which $U=0$ has probability at most $\kappa$,
by conditioning on $(P,\bar M)$.  Hence
$\Prob_i(\mathrm{Acc}_j\cap\{U=1\})\ge c/2$.

For the resampling, first fix the domains and projection choices, the path
through $D_i$, all matrices other than $M$, the omitted direction
$a\in K\setminus\{0\}$ for the level-$j$ test, and every bucket $h_v$ of
$M=\sum_{v\ne0}v\otimes h_v$ except $h_a$.  Denote this full background by
$\mathcal B$.  Fresh calls make $h_a$ and its resampled copy conditionally
independent with the same law.  The right level-$j$ query is constant as a
function of this bucket, because $M\bmod\langle a\rangle$ is constant.
Every successful left response therefore has the same lower outputs and
the same upper output modulo $\langle a\rangle$.

This assertion also holds for the representative upper output.  Indeed,
if $M'-M=a\otimes h$, linearity of the section gives
\[
 [M'-s(\bar M')]-[M-s(\bar M)]
        =a\otimes(h-s(\pi h)).
\]
Successful samples have the same $b_y$, so the corresponding folding
correction changes by a multiple of $a$.  Thus at fixed $\mathcal B$ the
useful successful samples have at most two possible values of
\eqref{eq:upper-representative}.  If their probabilities are $p_1,p_2$
(allowing one to be zero), two independent bucket draws yield equal defined
responses with probability at least
\[
 p_1^2+p_2^2\ge (p_1+p_2)^2/2.
\]
Averaging and applying Cauchy--Schwarz again proves the first bound.

The event just obtained is a function of $P,M,M'$ and the fixed partial
tables.  It records neither the selected descendant below $D_j$ nor the
projection metadata.  Apply the modified-$i$ to unrestricted-$j$ comparison
from Lemmas~\ref{lem:cut-tv} and~\ref{lem:sparse-tv}, including the additional
bucket call.  Own buckets at $D_j$, pending calls used in its ancestors,
and descendant matrices are separate independent families in the reference
experiment.  Consequently $M$ is uniform independently of $P$: independent
uniform buckets map surjectively to $K\otimes H$.  Given all old buckets,
$h=h_a+h_a'$ is uniform in $H$ and independent of those buckets.  Hence
$M'=M+a\otimes h$ is precisely the uniform rank-one walk, and its image in
$H/B$ is the same walk on the quotient.  The claimed loss gives
$c^2/8-c^2/16=\theta$.
\end{proof}

\subsection{From row advice to a predictable upper response}

Write $\Prob_j^0$ for the unrestricted cut law in the last assertion of
Lemma~\ref{lem:useful-equality}, with independent auxiliary choices appended
as needed.  Put $\eta=\theta/2$.  If the average conditional equality given $P$
is at least $\theta=2\eta$, the set of $P$ for which it is at least
$\eta$ has probability at least $(\theta-\eta)/(1-\eta)\ge\eta$.  Use the constants
$0<\alpha\le1$, $r\ge1$ in Lemma~\ref{lem:many-fibers} at equality threshold
$\eta$, obtained from Theorem~\ref{thm:inverse} at threshold $\eta/2$.
Assume that $\ell$ meets their largeness requirement.

Sample $A\in\Hom(K,\F^r)$ uniformly and independently.  For any
$(P,A,C)$, where $C\in\F^r\otimes(H/B)$, call the advice \emph{good} if
the row fiber $A\bar M=C$ has a nonempty slice obtained by specifying at
most $r$ columns on which one defined value of $\widetilde f_P$ has density
at least $\alpha$ in uniform slice measure.  For each good advice choose,
using fixed deterministic orders, column functionals
$q_1,\ldots,q_k\in(H/B)^*$, values $v_1,\ldots,v_k\in K$, and a response
\[
 t_*=(u,(y_C)_{C\in\mathcal D_i(D_j)}),\qquad k\le r,
\]
witnessing this condition.  These choices depend only on the advice; they
are not selected using the actual position within the column slice.
The probability of good advice $(P,A,A\bar M)$ under $\Prob_j^0$ is at least
\begin{equation}\label{eq:upper-g}
 g=\frac{\eta^2}{4}\,2^{-r\ell}.
\end{equation}
Indeed, the set of qualifying $P$ has mass at least $\eta$, and for each
such $P$ Lemma~\ref{lem:many-fibers} gives good-advice probability at least
$(\eta/4)2^{-r\ell}$.  Row maps of deficient rank are allowed throughout.

Pull the chosen $q_a$ back to $H^*$ and put
$Z=\Span\{q_1,\ldots,q_k\}\subseteq\Ann(B)$.  The chosen response occurs
on a nonempty set of legal representative queries, so the preceding
joint-legality argument defines its functional $b_y\in B^*$.
Define $z_0\in H^*$ by
\begin{equation}\label{eq:upper-z0}
 z_0(b+s(\bar h))=b_y(b)
       \quad(b\in B,\ \bar h\in H/B).
\end{equation}
This uses the direct-sum decomposition $H=B\oplus s(H/B)$.  At every match
to the selected response the restored upper output equals
\begin{equation}\label{eq:upper-restored}
 Mz_0+u.
\end{equation}
Every $z\in z_0+Z$ satisfies
\begin{equation}\label{eq:upper-valid}
 z(1)=1,\qquad
 \bigl(z(s_{C,t}s_{C,u})\bigr)_{t,u}
       =y_Cy_C^{\mathsf t}\quad(C\in\mathcal D_i(D_j)).
\end{equation}
Thus the whole coset is normalized and has tested rank at most one at
each of these descendants.  Dependent generators in $B$ cause no ambiguity;
legality, rather than a choice of their representation, defined $b_y$.

Define two events on $(P,A,M)$:
\begin{align*}
 J&=\{\text{advice is good and }Mq_a=v_a\text{ for }1\le a\le k\},\\
 I&=J\cap\{\widetilde f_P(\bar M)=t_*\}.
\end{align*}
Both events belong to $\sigma(P,A,\bar M)$, since every $q_a$ annihilates
$B$ and the witness uses only $(P,A,A\bar M)$.  In particular, $I$ asks
for a representative-table match; it does not ask for a restored right
answer to equal \eqref{eq:upper-restored}.
For uniform $M$ in a fixed nonempty row fiber, imposing $k$ columns costs
at most $k\ell\le r\ell$ binary equations.  Therefore, with
\begin{equation}\label{eq:upper-q0}
 h_0=2^{-r\ell},\qquad q_0=gh_0,
\end{equation}
the reference law satisfies
\begin{equation}\label{eq:upper-IJ}
 \Prob_j^0(J)\ge q_0,\qquad
 \Prob_j^0(I)\ge\alpha\Prob_j^0(J).
\end{equation}
The entire definition of goodness and the chosen witnesses is now fixed.
When evaluating $I,J$ under another law we use these same predicates.

\begin{lemma}[Predicting the upper response]
\label{lem:positive-difference}
Let $R_j^{(i)}\in K$ be the restored upper response of the level-$i$ right
table in the modified experiment.  Suppose the total-variation comparison
between $\Prob_i$ and $\Prob_j^0$, for the observations defining $I,J$, is
at most $v\le\alpha q_0/16$.  Then, with
\[
 \rho=\kappa\alpha/4,\qquad b=\kappa\alpha q_0/4,
\]
one has
\begin{equation}\label{eq:upper-difference}
 \Prob_i\bigl(J,\ R_j^{(i)}=Mz_0+u\bigr)-\rho\Prob_i(J)\ge b>0.
\end{equation}
All quantities $g,h_0,q_0,\rho,b$ depend only on $c,\alpha,r,\ell$.
\end{lemma}

\begin{proof}
On $I$ the mark $U$ is one.  Independence of $A$ and the measurability
just proved imply
\[
 \Prob_i(I\cap \mathrm{Acc}_i)
 =\E_i\bigl[\one_I\Prob_i(\mathrm{Acc}_i\mid P,\bar M,A)\bigr]
 \ge\kappa\Prob_i(I).
\]
Acceptance at level $i$ preserves the entire upper evaluation.  Thus on
$I\cap \mathrm{Acc}_i$ the right response is $Mz_0+u$.  Exact sharing of canonical
keys permits evaluating this right response directly on $O$.
Writing $h=\Prob_j^0(J)\ge q_0$, unconditional total variation and
\eqref{eq:upper-IJ} give
\[
 \Prob_i(I)\ge\alpha h-v,\qquad \Prob_i(J)\le h+v.
\]
The left side of \eqref{eq:upper-difference} is consequently at least
\[
 (\kappa\alpha-\rho)h-(\kappa+\rho)v
 \ge\frac{3\kappa\alpha}{4}q_0
      -\kappa(1+\alpha/4)\frac{\alpha q_0}{16}
 \ge b.
\]
This proves an unconditional inequality.  We will disintegrate it below,
without reusing the usefulness estimate after a finer conditioning.
\end{proof}

The affine witness now predicts the upper response with positive probability,
and every functional in its coset is normalized and passes the lower rank
test.  To obtain two local decoders, we need a bounded collection of candidate
functionals, computed from the right input, that contains the restriction of
some functional in this coset to the projected space.  The same collection
must work for every compatible left witness.

\subsection{The local inputs and the original hidden law}\label{sec:right-hidden-law}

Here and below the experiment is $\Prob_i$.  Write
$H'=H_{D_j}(O)$, $H_i=H_{D_i}(E)$, $H_i'=H_{D_i}(O)$, and
$S=M_{D_i}$.  Pullback identifies the primed spaces with subspaces of the
unprimed ones, preserving products.  Put $W=\ker A$ and split the upper
matrix into
\begin{equation}\label{eq:upper-hidden}
 M=M_{\rm known}+X,\qquad
 M_{\rm known}=\sum_{v\in K\setminus W}v\otimes h_v,\qquad
 X=\sum_{v\in W\setminus\{0\}}v\otimes h_v\in W\otimes H'.
\end{equation}
The known sum is determined without any bucket whose index lies in $W$.
In particular $AM=AM_{\rm known}$.

We specify the two inputs exactly.  Tree locations are retained in these
analytical inputs, even though the original partition keys omit unnecessary
sampling metadata.  The left information is
\begin{equation}\label{eq:upper-left-input}
 \mathcal L=\sigma\bigl(E,\text{path through }D_i,A,
       (M_D)_{D\ne D_j},(h_v)_{v\in K\setminus W}\bigr),
\end{equation}
where every function is represented on its original $E$-domain.  These
data supply $S$ in full, without supplying $a_i$, $O$, or the
projection choices as additional fields.  The right information is
\begin{equation}\label{eq:upper-right-input}
 \mathcal R=\sigma\bigl(O,\text{path through }D_i,A,a_i,
       (M_D)_{D\ne D_j,D_i},S\bmod\langle a_i\rangle,
       (h_v)_{v\in K\setminus W}\bigr),
\end{equation}
with functions represented on their $O$-domains.  The right is supplied
only the displayed quotient of $S$, with no additional missing-component
ingredients.  Neither side is supplied $X$.  The left can compute $P$ and $A\bar M$ from
\eqref{eq:upper-left-input}, so it can compute the good-advice predicate
and its chosen witness without knowing $M$.
Figure~\ref{fig:tree-cut} summarizes these information restrictions.

\begin{figure}[p]
\centering
\begin{tikzpicture}[
  >=Latex, font=\small,
  cut/.style={dashed,draw=black!50},
  treebox/.style={draw=blue!45!black,rounded corners=2pt,fill=blue!3,
    align=center,inner sep=5pt},
  inputbox/.style={draw=black!65,rounded corners=2pt,fill=yellow!4,
    align=left,text width=6.7cm,inner sep=7pt}]
\node[font=\small\bfseries] at (3.2,0) {Latent sampling structure};
\node[font=\small\bfseries,align=center] at (11.1,0.1) {Original sampler:\\separate calls at the same node};
\node at (3.2,-0.55) {ancestors};
\node[treebox] (upper) at (3.2,-1.25) {$D_j$: upper matrix $M$};
\draw[->] (3.2,-0.78)--(upper);
\draw[cut] (0,-1.25)--(1.0,-1.25);
\draw[cut] (5.4,-1.25)--(6.6,-1.25);
\node[align=left,font=\footnotesize] at (0.6,-1.78) {upper\\cut};
\node[align=center] (middle) at (3.2,-2.05) {$\vdots$};
\node[align=left,font=\footnotesize] at (5.45,-2.0) {$j>i$; not necessarily\\adjacent levels};
\node[treebox] (lower) at (3.2,-2.85) {$D_i$: lower matrix $S$};
\draw (upper)--(middle);
\draw (middle)--(lower);
\draw[cut] (0,-2.85)--(1.0,-2.85);
\draw[cut] (5.4,-2.85)--(6.6,-2.85);
\node[align=left,font=\footnotesize] at (0.6,-3.38) {lower\\cut};
\node (ordinary) at (1.65,-3.95) {ordinary $C$};
\node (special) at (4.7,-3.95) {$C_* = D_{i-1}$};
\draw (lower)--(ordinary);
\draw[thick,->] (lower)--(special);
\node[align=center,text width=6.5cm,font=\footnotesize] at (3.2,-4.65)
  {$C_*$ is the original sampler's latent child.\\
   The modified-$i$ law uses no path below $D_i$.};
\node[treebox,text width=5.2cm] (own) at (11.1,-1.45)
  {Own bucket at $D_p$\\$h_{p,v}\sim\nu_{D_p}$};
\node[treebox,text width=5.2cm] (call) at (11.1,-3.35)
  {Stopped ancestor call at $D_p$\\$h'\sim\nu_{D_p}$};
\draw[->] (11.1,-0.5)--node[right,font=\footnotesize]
  {fresh tape}(own.north);
\draw[->] (11.1,-2.4)--node[right,font=\footnotesize]
  {fresh tape}(call.north);
\node[align=center,text width=5.8cm,font=\footnotesize] at (11.1,-4.65)
  {Independent conditional on the domains and shared path;\\
   no call is reused in the other array.};
\node[inputbox,anchor=north west] (left) at (-0.2,-5.35) {
  \textbf{Left decoder: functions on $E$}\\[4pt]
  $E$, path through $D_i$, and $A$\\[3pt]
  Full $S$\\[3pt]
  $(M_D)_{D\ne D_j,D_i}$ on $E$\\[3pt]
  $(h_v)_{Av\ne0}$ on $E$\\[4pt]
  \footnotesize No $a_i$ or projection metadata is supplied.
};
\node[inputbox,anchor=north east] (right) at (14.5,-5.35) {
  \textbf{Right decoder: functions on $O$}\\[4pt]
  $O$, path through $D_i$, $A$, and $a_i$\\[3pt]
  Only $S\bmod\langle a_i\rangle$\\[3pt]
  $(M_D)_{D\ne D_j,D_i}$ on $O$\\[3pt]
  $(h_v)_{Av\ne0}$ on $O$\\[4pt]
  \footnotesize The missing component of $S$ is not supplied.
};
\node[fit=(left)(right),inner sep=0pt] (cards) {};
\node[align=center,text width=14cm,below=5mm of cards.south]
  {Neither decoder is supplied the hidden upper sum
   $X=\sum_{0\ne v\in\ker A}v\otimes h_v$.};
\end{tikzpicture}
\caption{Cuts and the information supplied to the decoders.
The selected levels satisfy $j>i$.  The tree depicts latent sampling
choices; the upper-cut comparison forgets the path below $D_j$.
The call panel shows the original sampler; at a modified cut these calls
are replaced by independent uniform draws in the cut space.
An own bucket and a stopped ancestor call at the same node use separate
tapes, independently conditional on the domains and shared path.
The cards summarize the analytical inputs
\eqref{eq:upper-left-input}--\eqref{eq:upper-right-input}, with $W=\ker A$.
The left receives full $S$ without $a_i$; the right receives $a_i$ and
only the quotient of $S$.  Both receive upper buckets with $Av\ne0$ on
their own domains.  These analytical records are distinct from the
original game's canonical partition keys.}
\label{fig:tree-cut}
\end{figure}


For clarity, the independence underlying these inputs is exact, not an
approximation.  Conditional on $O$, the projection choices, and the path
through $D_i$, give each actual node matrix its own random tape.  Give each
scalar bucket in a matrix a fresh recursive call tree, with independent
ordinary-child uniforms and independent factor calls at every product.
All calls that reach $D_i$ terminate in separate uniform draws in
$H_i'$.  In particular, even a call used for an ancestor of $D_j$ and a
bucket of $M$ that both visit $D_j$ use disjoint random tapes.
The product factorization of these tapes implies
\begin{equation}\label{eq:upper-factorization}
 (h_v)_{v\in W\setminus\{0\}}\mathrel{\perp\!\!\!\perp}
 \bigl((h_v)_{v\notin W},(M_D)_{D\ne D_j}\bigr)
 \quad\text{given the domains, path, and }A.
\end{equation}
The same statement remains true if the original questions $E$ and the
projection maps are additionally given: they determine the embeddings of
the spaces but do not change the intrinsic $O$-space sampler.  This also
explains why analytical disclosure of the full $S$ is harmless for this
factorization, although $S$ is not part of the right input.

Let $\mu_{\mathcal R}$ be the distribution on $W\otimes H'$ obtained by
sampling the hidden buckets in \eqref{eq:upper-hidden} with this original
modified-$i$ sampler.  It is a function of the right input alone.  Indeed,
the right knows $O$, the path down to the terminating cut, all spaces and
all sampler parameters.  Independent recursive calls above that cut can
be enumerated on those spaces without knowing $E$, any projection map,
or the missing part of $S$.  Formula~\eqref{eq:upper-factorization} shows
that $\mu_{\mathcal R}$ is also the actual conditional hidden law after
revealing all of the data in both inputs, including full $S$.

For each fixed right input, define
\begin{equation}\label{eq:upper-G}
 G_{\mathcal R}:W\otimes H'\longrightarrow K
\end{equation}
by assembling the level-$i$ right query with upper matrix
$M_{\rm known}+X$, consulting its fixed partition label, and restoring
the upper component.  This is a well-defined local function: the lower
matrix in this query is precisely $S\bmod\langle a_i\rangle$.
The original labeling has already been extended to unused legal keys, so
this definition is total for every $X$ in the ambient space.

Assume that every product count meets the absolute balance minimum of
Lemma~\ref{lem:bias}, and that
\begin{equation}\label{eq:upper-R}
 (7/8)^{R_j}\le\lambda.
\end{equation}
Every nonzero character of $W\otimes H'$ then has bias at most $\lambda$
under $\mu_{\mathcal R}$.  To check this directly, identify a character
with a nonzero linear map $\Phi:W\longrightarrow(H')^*$.  Some nonzero
$v\in W$ has $\Phi(v)\ne0$, so its expectation on the corresponding
bucket has absolute value at most $\lambda$ by Lemma~\ref{lem:bias}.
Independence factors the expectation over all hidden buckets, and all
remaining factors have absolute value at most one.  This argument makes
no claim about small bias after a slice, a successful match, or any later
conditioning.

\subsection{A bounded list for an arbitrary small-bias law}

We record the Fourier calculation separately.  It uses characters under
their actual sampling law rather than treating small bias as total-variation
closeness to uniform measure.  The short-list principle is related to
Goldreich--Levin heavy-character decoding~\cite{GoldreichLevin1989}.
We prove the needed statement for the present nonuniform, own-question law;
no efficient decoding algorithm is required.  If $V=W\otimes H'$, write
\[
 \chi_\Phi(X)=(-1)^{\langle\Phi,X\rangle},\qquad
 \Phi\in V^*=\Hom(W,(H')^*).
\]
For $\sigma\in K^*$ and a fixed function $G:V\longrightarrow K$, put
$g_\sigma(X)=(-1)^{\sigma(G(X))}$.

\begin{lemma}[One heavy list for all affine-slice witnesses]
\label{lem:heavy-list}
Let $K=\F^\ell$, let $W\subseteq K$ have codimension at most $r$, and
let $H'$ be a finite vector space.  Let $\mu$ be a probability law on
$V=W\otimes H'$ with $|\E_\mu\chi_\Phi|\le\lambda$ for every nonzero
$\Phi\in V^*$.  Fix $G:V\longrightarrow K$, $0<\rho\le1$, and set
\[
 h_0=2^{-r\ell},\qquad \tau=\rho h_0/4.
\]
Suppose
\begin{equation}\label{eq:upper-Fourier-conditions}
 2^{-(\ell-r)}<\rho/8,
 \qquad 0\le\lambda\le\min\{h_0/2,\tau^2/2\}.
\end{equation}
For each $\sigma\in K^*$ define the list
\begin{equation}\label{eq:upper-heavy-definition}
 \mathcal H_\sigma
 =\{\Phi\in V^*:|\E_\mu[g_\sigma\chi_\Phi]|\ge\tau\}.
\end{equation}
Each list has size at most $2/\tau^2$.  Moreover, let
$Q\subseteq(H')^*$ have dimension at most $r$, let
$t_Q\in\Hom(Q,W)$ specify the nonempty affine slice
$\mathcal C=\{X:X|_Q=t_Q\}$, and let
$z'\in(H')^*$ and $k_0\in K$.  If
\begin{equation}\label{eq:upper-slice-equality}
 \Prob_\mu(G(X)=Xz'+k_0\mid X\in\mathcal C)\ge\rho,
\end{equation}
then some $\sigma$ with $\sigma|_W\ne0$ has a list element in
\begin{equation}\label{eq:upper-heavy-coset}
 (\sigma|_W)\otimes z'+W^*\otimes Q.
\end{equation}
The lists in \eqref{eq:upper-heavy-definition} depend on $\mu,G,W,H'$ and
$\tau$, and not on the witness $Q,t_Q,z',k_0$.
\end{lemma}

\begin{proof}
Put $w=\dim W$ and $q=\dim Q$.  The map
$V\longrightarrow\Hom(Q,W)$ given by $X\mapsto X|_Q$ is surjective:
choose a basis of $Q$, extend it to a basis of $(H')^*$, and extend any
specified linear map to all of $(H')^*$.  Thus the linear equations for
$\mathcal C$ have rank $wq\le\ell r$.  Their dual space is
$L=W^*\otimes Q\subseteq V^*$.  Choose any $X_0\in\mathcal C$.  The
exact character identity for the affine slice is
\begin{equation}\label{eq:upper-slice-expansion}
 \one_{\mathcal C}(X)
   =\frac1{|L|}\sum_{\Psi\in L}\chi_\Psi(X_0)\chi_\Psi(X).
\end{equation}
Indeed, the sum is one when $X-X_0$ annihilates $L$, and zero otherwise.
Consequently
\begin{equation}\label{eq:upper-slice-mass}
 \mu(\mathcal C)\ge2^{-wq}-\lambda\ge h_0-\lambda\ge h_0/2>0.
\end{equation}
This proves the needed slice mass using unconditional bias.

Expand the equality in \eqref{eq:upper-slice-equality} over $K^*$:
\[
 \Prob_\mu(G(X)=Xz'+k_0\mid\mathcal C)
 =2^{-\ell}\sum_{\sigma\in K^*}(-1)^{\sigma(k_0)}
    \E_\mu[g_\sigma(X)\chi_{(\sigma|_W)\otimes z'}(X)
                    \mid\mathcal C].
\]
The fraction of $\sigma$ annihilating $W$ is $2^{-w}$, which is at most
$2^{-(\ell-r)}<\rho/8$.  Each term has absolute value at most one, so
some remaining $\sigma$ has
\[
 \left|\E_\mu[g_\sigma\chi_{(\sigma|_W)\otimes z'}\mid\mathcal C]\right|
 \ge\rho/2.
\]
Multiply by \eqref{eq:upper-slice-mass} and then insert
\eqref{eq:upper-slice-expansion}.  The coefficients in that expansion
have total absolute value one.  It follows that for some $\Psi\in L$,
\[
 \left|\E_\mu[g_\sigma
       \chi_{(\sigma|_W)\otimes z'+\Psi}]\right|
 \ge(\rho/2)(h_0/2)=\tau.
\]
This gives \eqref{eq:upper-heavy-coset}.  Notice that the unknown target
offset contributes only a sign in the equality expansion.

It remains to bound the lists under the possibly nonuniform law $\mu$.
Take $m$ distinct elements $\Phi_1,\ldots,\Phi_m$ of one list and choose
signs $\epsilon_a\in\{-1,1\}$ aligning their correlations with
$g_\sigma$.  Cauchy--Schwarz gives
\begin{align*}
 m^2\tau^2
 &\le\left(\E_\mu\left[g_\sigma
               \sum_{a=1}^m\epsilon_a\chi_{\Phi_a}\right]\right)^2\\
 &\le\E_\mu\left[\left(\sum_{a=1}^m
                   \epsilon_a\chi_{\Phi_a}\right)^2\right]
 \le m+m(m-1)\lambda\le m+m^2\lambda.
\end{align*}
For distinct indices the product character is nonzero, which justifies
the off-diagonal bound even after sign alignment.  If $m>0$, the hypothesis
$\lambda\le\tau^2/2$ yields $m\le2/\tau^2$; the empty case is immediate.
Only the selected characters were summed, so no ambient-dimension factor
appears.  All correlations defining the lists are unconditional under the
same law $\mu$.
\end{proof}

\subsection{Constructing the two decoders}

\begin{proposition}[Upper decoding with a lower rank test]
\label{prop:upper}
Suppose a legal labeling has acceptance at least $\eps$, $d\eps\ge4$,
and the cut comparisons have the errors required in
\eqref{eq:upper-c}, Lemma~\ref{lem:useful-equality}, and
Lemma~\ref{lem:positive-difference}.  Choose upper inverse constants
$\alpha,r$ and $\ell=\ell_j$ as above, and assume
\eqref{eq:upper-Fourier-conditions} and \eqref{eq:upper-R}, with the
absolute balance minimum at every level.  For the pair $i<j$ supplied by
\eqref{eq:upper-pair}, there are randomized strategies, measurable on
\eqref{eq:upper-left-input} and \eqref{eq:upper-right-input} respectively,
returning bilinear forms $F_L$ on $H_i$ and $F_R$ on $H_i'$ such that
\begin{equation}\label{eq:upper-gamma}
 \Prob_i\bigl(F_R=F_L|_{H_i'\times H_i'}\bigr)\ge
 \gamma_j:=b\,2^{-\ell}\frac{\tau^2}{2}\,2^{-r}>0.
\end{equation}
Every left output has the form
\[
 F_L(x,y)=z(xy),\quad z\in(H_i^2)^*,\quad z(1)=1,
 \qquad \rank\bigl(F_L(S_t,S_u)\bigr)_{t,u}\le1.
\]
The constant $\gamma_j$ depends only on $c,\alpha,r,\ell_j$; it is
independent of the lower row count $\ell_i$, the branching numbers, the
source alphabet and instance size, and the bias accuracies chosen for
lower levels.  The left strategy never receives $a_i$, and the right
strategy uses $S$ only through its quotient by $\langle a_i\rangle$.
\end{proposition}

\begin{proof}
On good advice the left samples $z$ uniformly from $z_0+Z$, using fresh
private randomness, and returns $F_L(x,y)=z(xy)$ for $x,y\in H_i$.
The products belong to $H$ because $D_i$ is a proper descendant of $D_j$.
Equation~\eqref{eq:upper-valid} proves normalization and the rank bound
on $S$.  On bad advice it evaluates at a fixed legal assignment on the
original subtree, obtaining the same two validity properties.  These rules
also define valid defaults on any other well-typed left input: use the
stated construction whenever its witness is valid, and point evaluation
otherwise.  The rule involves no omitted direction.

For the right decoder, use $\mu=\mu_{\mathcal R}$ and
$G=G_{\mathcal R}$ from \eqref{eq:upper-G}, and compute all lists
\eqref{eq:upper-heavy-definition}.  Choose $\sigma\in K^*$ uniformly.
If $\sigma|_W\ne0$ and its list is nonempty, choose an element $\Phi$
uniformly from that list.  Choose $h_\sigma\in W$ deterministically from
$(A,\sigma)$ with $\sigma(h_\sigma)=1$, for example the first such vector
in a fixed order, and put $z_R=\Phi(h_\sigma)\in(H')^*$.  Return
\begin{equation}\label{eq:upper-right-form}
 F_R(x,y)=z_R(xy),\qquad x,y\in H_i'.
\end{equation}
The inclusion $(H_i')^2\subseteq H'$ makes this a bilinear form.  In the
remaining cases return the zero bilinear form.  These defaults make the
algorithm total; normalization is required only of the left output.
No chosen witness from the left enters the right rule.

We now prove agreement for these particular algorithms.  For this analysis
only, expose $E,O$, all projection maps, the path through $D_i$, $A,a_i$,
all nonhidden upper buckets, and all other matrices including full $S$.
Call the resulting data $\mathcal D$.  Do not expose $X$, any hidden
bucket ingredients, the slice event $J$, the representative-match event
$I$, a successful test, or either decoder's final random choices.
This exposure determines $P$, $M_{\rm known}$ and $A\bar M$, since
$AX=0$.  Hence it determines goodness, $Z,z_0,u$ and the column witness.
By \eqref{eq:upper-factorization}, the conditional distribution of $X$
given $\mathcal D$ remains exactly $\mu_{\mathcal R}$, with the bias
bound already proved.

On good data let $Q=Z|_{H'}\subseteq(H')^*$ and $z'=z_0|_{H'}$.
The event $J$ either is impossible, or is an affine slice
\[
 \mathcal C_{\mathcal D}=
 \{X\in W\otimes H':X(q_a|_{H'})=v_a-M_{\rm known}q_a
                         \text{ for every }a\}.
\]
In the latter case these equations define a well-defined linear
specification on $Q$, of dimension at most $r$.  This also handles a
collapse of some columns on restriction to $H'$.  The target in
\eqref{eq:upper-difference} is $Xz'+k_0$, where
$k_0=M_{\rm known}z_0+u$ is fixed by $\mathcal D$.

Set $p_{\mathcal D}=\mu_{\mathcal R}(\mathcal C_{\mathcal D})$ on good
compatible data, and zero otherwise.  When $p_{\mathcal D}>0$, let
\[
 q_{\mathcal D}=
 \Prob_i\bigl(G_{\mathcal R}(X)=Xz'+k_0
                         \mid\mathcal D,\ X\in\mathcal C_{\mathcal D}\bigr),
\]
and set $q_{\mathcal D}=0$ otherwise.  Disintegrating the already proved
unconditional inequality \eqref{eq:upper-difference} gives
\begin{equation}\label{eq:upper-favorable-mass}
 \E_i[p_{\mathcal D}(q_{\mathcal D}-\rho)]\ge b.
\end{equation}
Its integrand is nonpositive unless the data are good and compatible,
$p_{\mathcal D}>0$, and $q_{\mathcal D}\ge\rho$.  The integrand is at
most one everywhere.  The set of data with those properties consequently
has probability at least $b$.  This is the sole use of the finer exposure:
we have not asserted that usefulness persists upon conditioning on it.

Fix any such favorable $\mathcal D$.  Lemma~\ref{lem:heavy-list} supplies
at least one nonannihilating $\sigma$ and one
\[
 \Phi\in\mathcal H_\sigma\cap
       \bigl((\sigma|_W)\otimes z'+W^*\otimes Q\bigr).
\]
The right chooses such a pair with probability at least
$2^{-\ell}\tau^2/2$.  Every such choice, with the already fixed
$h_\sigma$, has
\[
 z_R=\Phi(h_\sigma)\in z'+Q.
\]
The restriction map $Z\longrightarrow Q$ is surjective.  Uniform measure
on $z_0+Z$ therefore pushes to uniform measure on $z'+Q$, so the fresh
left choice has restriction equal to this particular $z_R$ with probability
$2^{-\dim Q}\ge2^{-r}$.  It is independent of the right's final choices
conditional on $\mathcal D$.  Equality of these functionals on $H'$
implies agreement of their product forms on $H_i'\times H_i'$.
Averaging over the favorable data proves \eqref{eq:upper-gamma}.

For a fixed right input, the law $\mu_{\mathcal R}$, function
$G_{\mathcal R}$, lists $\mathcal H_\sigma$, and contraction vectors
$h_\sigma$ are identical for all compatible left inputs and witnesses.
Different favorable exposures may certify different elements of these same
lists.  The preceding hit probability holds separately for each exposure,
so averaging introduces no factor for the number of witnesses.

Finally, $\eta,\alpha,r$ were fixed from $c$, while $g,q_0,b,\rho,\tau$
were fixed from these constants and $\ell_j$.  Conditions
\eqref{eq:upper-Fourier-conditions} require only an upper row-count
lower bound and an upper bias tolerance.  Condition~\eqref{eq:upper-R}
is met by choosing $R_j$ large enough; below $j$ it uses only the common
absolute balance minimum.  The inverse statement applies to $H/B$ even
in small dimension by Lemma~\ref{lem:padding}.  Thus none of these choices
depends on the later lower row count or source and branching dimensions.
\end{proof}

The algorithms in Proposition~\ref{prop:upper} are now fixed local rules.
In particular, the right rule always computes its heavy lists from the
original law $\mu_{\mathcal R}$ just defined.  A later conditioning may
change the physical law of the hidden buckets; it does not change that
algorithm.  The subsequent rank and repetition arguments use these same
rules with their prescribed inputs.
